% Lösungen Einheit 2 von 5 – Lineare Funktion f(x) = m·x + n (zusammenbau v0.1)
\einheitenkopf{Einheit 2 von 5 · Lineare Funktion f(x) = m·x + n}

\erg{12}{a) $m = 3$, $n = -5$}
\erg{13}{a) steigt (m ist positiv)}
\erg{14}{a) 3 nach oben, also $m = 3$}
%% TODO Lösung der Erklärzeile 15f) fehlt in der Bank
\erg{15}{a) $-5$, $-2$, $1$, $4$ \quad b) $-1$, $1$, $3$, $5$ \quad c) $5$, $4$, $3$, $2$ \quad d) $-7$, $-3$, $1$, $5$ \quad e) $5$, $3$, $1$, $-1$ \quad f) \ldots \quad g) Punkte $(-2|-2)$, $(-1|1)$, $(0|4)$, $(1|7)$; sie liegen auf einer Geraden \quad h) zum Beispiel $(0|2)$ und $(1|6)$; Gerade durch beide Punkte \quad i) Start $(0|5)$, 1 nach rechts, 3 nach oben: Gerade durch $(0|5)$ und $(1|8)$ \quad j) Start $(0|4)$, 1 nach rechts, 3 nach unten: Gerade durch $(0|4)$ und $(1|1)$ \quad k) Start $(0|3)$, 2 nach rechts, 1 nach oben: Gerade durch $(0|3)$ und $(2|4)$ \quad l) Start $(0|-4)$, 1 nach rechts, 2 nach oben: Gerade durch $(0|-4)$ und $(1|-2)$ \quad m) Start $(0|3)$, 1 nach rechts, 4 nach unten: Gerade durch $(0|3)$ und $(1|-1)$}
%% TODO Lösung der Erklärzeile 16f) fehlt in der Bank
\erg{16}{a) $n = 3$ \quad b) $n = 4$ \quad c) $n = -2$ \quad d) $n = -1$ \quad e) $n = -4$ \quad f) \ldots \quad g) $m = 2$ (1 nach rechts, 2 nach oben) \quad h) $m = -2$ (1 nach rechts, 2 nach unten) \quad i) $m = \frac{1}{2}$ (2 nach rechts, 1 nach oben) \quad j) $f(x) = 3x - 4$ \quad k) f: $y = -3x + 3$, g: $y = 2x + 3$, h: $y = -3$}
\erg{17}{a) f und h, beide haben $m = 4$}
\erg{18}{a) Paul hat „nach rechts durch nach oben“ gerechnet statt „nach oben durch nach rechts“; richtig ist $m = 3$. \quad b) g, denn 5 ist größer als 3: je Schritt nach rechts geht g weiter nach oben. \quad c) $m = -1$, $n = 3$, also $f(x) = -x + 3$}
